\section{Análise no Tempo Discreto}

A análise no tempo discreto foi realizada com o objetivo de avaliar como a dinâmica contínua do sistema pode ser representada por amostras tomadas em instantes definidos de tempo.
Essa etapa é relevante porque, em aplicações industriais, mesmo que o processo físico seja contínuo, as medições e as ações de controle são normalmente executadas por sistemas digitais.
Assim, variáveis como temperatura, fração molar e vazão são lidas em intervalos de amostragem e utilizadas por controladores digitais, sistemas supervisórios ou algoritmos de controle avançado.

Para as análises apresentadas a seguir, foi considerada a função de transferência que relaciona a vazão de reciclo $F_R$ com a fração molar do componente $A$ no primeiro reator, $x_{A1}$.
% Essa relação é representada por:

% \begin{equation}
% G_{x_{A1},F_R}(s)=\frac{X_{A1}(s)}{F_R(s)}
% \end{equation}

% Essa escolha é coerente com o comportamento dinâmico observado nas etapas anteriores, pois a vazão de reciclo possui influência direta sobre a composição do primeiro reator.
%Além disso, a resposta de $x_{A1}$ a uma perturbação em $F_R$ apresenta comportamento transitório relevante, incluindo resposta inversa, o que torna essa combinação entre entrada e saída adequada para discutir a escolha do período de amostragem e a estabilidade no plano discreto.

\subsection{Amostragem da resposta temporal}

A escolha do período de amostragem ($T_s$) é uma etapa fundamental na representação discreta de um sistema dinâmico.
Quando $T_s$ é muito elevado, a resposta passa a ser descrita por poucos pontos, podendo haver perda de informações importantes sobre o comportamento transitório.
Além disso, a frequência de amostragem pode tornar-se insuficiente para representar adequadamente as componentes de alta frequência do sinal, ocasionando o fenômeno de \textit{aliasing}, que resulta em uma representação incorreta da dinâmica do sistema.
Por outro lado, um período de amostragem excessivamente pequeno eleva a quantidade de dados a serem processados e o esforço computacional, sem resultar em ganhos significativos na representação da dinâmica do sistema.

Dessa forma, o período de amostragem foi definido a partir da análise do espectro da resposta temporal, obtido por meio da Transformada Rápida de Fourier (\textit{Fast Fourier Transform} — FFT).
O espectro da resposta de $x_{A1}$ a um degrau unitário aplicado na entrada $F_R$, utilizado nessa análise, é apresentado no Apêndice~\ref{sec:apendice_fft}.
A partir desse espectro, estimou-se uma frequência máxima relevante de aproximadamente $\omega_{\max}=150\ \mathrm{rad/h}$.
Com base no Teorema de Amostragem de Shannon, adotou-se uma frequência angular de amostragem de $\omega_{\mathrm{amost}}=300\ \mathrm{rad/h}$, correspondente a um período de amostragem de $T_s\approx0{,}0209\ \mathrm{h}$, utilizado para a discretização do sinal.

A Figura~\ref{fig:amostragem_xa1_fr} apresenta a comparação entre o sinal contínuo e os pontos amostrados.
Observa-se que os pontos amostrados acompanham adequadamente a forma da resposta contínua.
Isso indica que o período de amostragem escolhido foi adequado para preservar as principais características do comportamento transitório do sistema, incluindo a resposta inversa e o tempo de acomodação.

\begin{figure}[htbp]
  \centering
  \includegraphics[width=0.75\textwidth]{figuras/amostragem.png}
  \caption{Resposta contínua e amostrada de $x_{A1}$.}
  \label{fig:amostragem_xa1_fr}
\end{figure}

\subsection{Mapeamento dos polos para o plano discreto}

Além da análise da resposta discretizada, foi avaliada a estabilidade do sistema no domínio discreto.
Para isso, os polos da função de transferência contínua foram mapeados para o plano $z$ por meio da relação:
\begin{equation}
  z_i = e^{s_iT_s},
\end{equation}
em que $s_i$ representa um polo no plano contínuo e $z_i$ representa o polo correspondente no plano discreto.
No tempo contínuo, a estabilidade de um sistema linear está associada à localização dos polos no semiplano esquerdo do plano $s$.
No tempo discreto, por sua vez, a estabilidade é verificada pela localização dos polos no plano $z$, sendo necessário que todos estejam no interior do círculo unitário, isto é:
\begin{equation}
  |z_i| < 1.
\end{equation}

Aplicando o mapeamento com $T_s = 0{,}0209 \ \mathrm{h}$, foram obtidos os polos discretos correspondentes aos polos do modelo contínuo.
Os valores numéricos são apresentados no Apêndice~\ref{sec:apendice_polos}.
A Figura~\ref{fig:pzplot_discreto} mostra a localização dos polos discretos no plano $z$.
A região delimitada pelo círculo unitário corresponde à região de estabilidade para sistemas discretos.
Observa-se que todos os polos estão localizados no interior desse círculo, indicando que a função de transferência discretizada satisfaz o critério de estabilidade no domínio discreto.
Além disso, nota-se que o polo mais próximo do círculo unitário é o polo dominante, estando associado ao modo mais lento da resposta, enquanto os demais polos, por estarem mais afastados, representam modos de decaimento mais rápido.

\begin{figure}[htbp]
  \centering
  \includegraphics[width=0.45\textwidth]{figuras/pzplot.png}
  \caption{Localização dos polos discretos da função de transferência entre $F_R$ e $x_{A1}$ em relação ao círculo unitário no plano $z$.}
  \label{fig:pzplot_discreto}
\end{figure}
